% Integration material requested by the author for Claude's v11.
% This is a section fragment, not a manuscript version.
% Base inspected: prime_corrected_clock_v9.tex, commit 2d8e702855192f2d959f66598a0289c31c4bfb28.
% No zero ordinates or zero-dependent parameters enter the construction.
% New labels use reg:; new local symbols have a reg superscript where needed.

\section{A regularized clock with a fixed prime-cutoff exponent}
\label{reg:section}

Fix $0<c<1/2$ independently of $T$ and put $X=T^c$.
The construction in this section uses only $\NO$, the Euler phase $S_X$,
and deterministic restrictions on the slope of a counting model. It does
not evaluate $S$ or use ordinates in constructing the points. The estimates
are unconditional, with constants depending on the fixed exponent $c$.
They concern a new regularized clock; they do not assert bounded transport
for the original one-step clock $\tau_k-S_X(\tau_k)/\lL(\tau_k)$.

\subsection{A finite convex construction}
Set
\begin{equation}
\begin{split}
\mathcal J_T^{\mathrm{reg}}&=[T/2,3T],\qquad
\mathsf a_T^{\mathrm{reg}}=\tfrac12\lL(T),\qquad
\mathsf b_T^{\mathrm{reg}}=2\lL(T),\\
\mathsf m_T^{\mathrm{reg}}&=\left\lceil\tfrac52T\lL(T)\right\rceil,\qquad
\mathsf h_T^{\mathrm{reg}}=\frac{5T}{2\mathsf m_T^{\mathrm{reg}}},\qquad
\mathsf t_{T,j}^{\mathrm{reg}}=T/2+j\mathsf h_T^{\mathrm{reg}}quad
(0\le j\le\mathsf m_T^{\mathrm{reg}}).
\end{split}
\label{reg:mesh}
\end{equation}
Thus $\mathsf h_T^{\mathrm{reg}}\le1/\lL(T)$.
For a vector $\mathbf v=(v_0,\ldots,v_{\mathsf m_T^{\mathrm{reg}}})$,
let $\mathsf p_{\mathbf v,T}^{\mathrm{reg}}$ be its continuous linear
interpolant on this mesh. Let $\mathcal C_T^{\mathrm{reg}}$ be the set
of these interpolants satisfying
\begin{equation}
\mathsf a_T^{\mathrm{reg}}\mathsf h_T^{\mathrm{reg}}
\le v_{j+1}-v_j\le
\mathsf b_T^{\mathrm{reg}}\mathsf h_T^{\mathrm{reg}}
\quad(0\le j<\mathsf m_T^{\mathrm{reg}}).
\label{reg:constraints}
\end{equation}
Define
\begin{equation}
\begin{split}
\mathscr E_{T,c}^{\mathrm{reg}}(\mathbf v)
&=\int_{\mathcal J_T^{\mathrm{reg}}}
\left|\mathsf p_{\mathbf v,T}^{\mathrm{reg}}(t)-F_X(t)\right|^4\,dt,\\
\widehat F_{T,c}
&=\underset{P\in\mathcal C_T^{\mathrm{reg}}}{\operatorname{argmin}}
\int_{\mathcal J_T^{\mathrm{reg}}}|P(t)-F_X(t)|^4\,dt.
\end{split}
\label{reg:projection}
\end{equation}
The admissible set is nonempty, closed and convex in a finite-dimensional
space. The objective is coercive and strictly convex: distinct nodal vectors
give interpolants differing on a set of positive measure. Consequently the
minimizer exists and is unique.

For indices whose target levels lie in its range, define
\begin{equation}
\widehat F_{T,c}(\widehat x_{T,c,k})=k-\tfrac12=\NO(\tau_k).
\label{reg:clock}
\end{equation}
There is at most one solution, since the lower slope is positive. If the
target lies between the consecutive minimizing values $\widehat v_j$ and
$\widehat v_{j+1}$, it is given directly by
\begin{equation}
\widehat x_{T,c,k}=\mathsf t_{T,j}^{\mathrm{reg}}+
\mathsf h_T^{\mathrm{reg}}
\frac{k-\frac12-\widehat v_j}{\widehat v_{j+1}-\widehat v_j}.
\label{reg:inverse}
\end{equation}
No branch choice or Newton recurrence occurs. The interval, mesh and slope
restrictions are prescribed analytically, before any comparison with zeros.

\subsection{Counting functions with controlled slopes}
The following proof uses the true count only to exhibit a comparison
function in the admissible class. This comparison function is not an input
to \eqref{reg:projection}.

\begin{lemma}\label{reg:count_comparator}
For all sufficiently large $T$ there exists $P_N\in\mathcal C_T^{\mathrm{reg}}$
such that
\begin{equation}
\int_{\mathcal J_T^{\mathrm{reg}}}|P_N-N|^4\,dt\ll T.
\label{reg:comparator_bound}
\end{equation}
The implied constant is absolute.
\end{lemma}

\begin{proof}
For this proof only, abbreviate $a=\mathsf a_T^{\mathrm{reg}}$ and
$b=\mathsf b_T^{\mathrm{reg}}$ and define
\[
\mathsf H_T^{\mathrm{reg}}(t)=
\sup_{s\in\mathcal J_T^{\mathrm{reg}}}
\{N(s)+\min(a(t-s),b(t-s))\}.
\]
Every function inside the supremum has increments between $a(t_2-t_1)$
and $b(t_2-t_1)$ for $t_2\ge t_1$. The supremum has the same property;
it is finite, continuous, and majorizes $N$ on $\mathcal J_T^{\mathrm{reg}}$.
Its excess over $N(t)$ is the maximum of zero and
\[
\sup_{\substack{h\ge0\\t-h\in\mathcal J_T^{\mathrm{reg}}}}
\{ah-N(t)+N(t-h)\},\qquad
\sup_{\substack{h\ge0\\t+h\in\mathcal J_T^{\mathrm{reg}}}}
\{N(t+h)-N(t)-bh\}.
\]
For $T$ sufficiently large,
$3\lL(T)/4\le\NO'(u)\le5\lL(T)/4$ on $[T/3,4T]$.
For the left increment use $j=\lfloor h\lL(T)\rfloor$;
for the right increment use $j=\lceil h\lL(T)\rceil$.
The monotonicity of $N$ and $N=\NO+S+O(1/T)$ give
\begin{equation}
0\le\mathsf H_T^{\mathrm{reg}}(t)-N(t)
\le4+\sup_{j,\,\pm}
\left(\left|\mathsf D_{T,j}^{\pm,\mathrm{reg}}(t)\right|-j/4\right)_+,
\label{reg:excess}
\end{equation}
where $j\ge1$ and
$\mathsf D_{T,j}^{\pm,\mathrm{reg}}(t)=S(t\pm j/\lL(T))-S(t)$.
Each increment is included only when its endpoints are in $[T/3,4T]$.
Indeed the deterministic drifts are at most $-j/4+1/2$ on the left
and $-3j/4+2$ on the right, before the $O(1/T)$ remainder.
One grid extension at an endpoint of $\mathcal J_T^{\mathrm{reg}}$
stays inside $[T/3,4T]$. The case $j=0$ on the left is at most $1/2$.

Theorem~\ref{thm:known_fujii_difference} with the fixed order $k=4$
gives, for shifts at most $1$,
\begin{equation}
\int_{\Omega_{T,j}^{\pm}}
|\mathsf D_{T,j}^{\pm,\mathrm{reg}}(t)|^8\,dt
\le C_{\mathrm D}T\log^4(2+j).
\label{reg:increment_eighth}
\end{equation}
Here $\Omega_{T,j}^{\pm}$ is the subset just described. Its bounded
dyadic covering has constants independent of $j,T$. For shifts greater
than $1$, $j>\lL(T)$, and the same estimate follows from
$|z-w|^8\le2^7(|z|^8+|w|^8)$ and
Theorem~\ref{thm:known_selberg_moments} with the fixed order $4$:
the integral is $O(T(\log\log T)^4)$, which is bounded by the right
side of \eqref{reg:increment_eighth}. Both shifted arguments stay in
$[T/3,4T]$. No estimate uniform in a growing moment order is used.

For $z\ge0$ and $j\ge1$,
$(z-j/4)_+^4\le4^4z^8/j^4$.
The fourth power of a supremum of nonnegative terms is bounded by the
sum of their fourth powers. Therefore
\[
\int_{\mathcal J_T^{\mathrm{reg}}}
|\mathsf H_T^{\mathrm{reg}}-N|^4\,dt
\le C_{\mathrm H}T,
\quad
C_{\mathrm H}=5120+4096C_{\mathrm D}
\sum_{j\ge1}\frac{\log^4(2+j)}{j^4}<\infty.
\]
Linear interpolation of $\mathsf H_T^{\mathrm{reg}}$ on the prescribed
mesh preserves the slope restrictions. It changes the function by at
most $b\mathsf h_T^{\mathrm{reg}}\le2$. Its interpolant is $P_N$, and
one can take
\begin{equation}
C_{\mathrm{cmp}}=8C_{\mathrm H}+320
\label{reg:comparator_constant}
\end{equation}
in \eqref{reg:comparator_bound}.
\end{proof}

\begin{lemma}\label{reg:continuous_fourth}
For every fixed $0<c<1/2$ there is $C_{\mathrm{pr}}(c)$ such that
\begin{equation}
\int_{\mathcal J_T^{\mathrm{reg}}}|F_{T^c}-N|^4\,dt
\le C_{\mathrm{pr}}(c)T
\label{reg:prime_fourth}
\end{equation}
for all sufficiently large $T$, and
\begin{equation}
\int_{\mathcal J_T^{\mathrm{reg}}}|\widehat F_{T,c}-N|^4\,dt
\le C_{\mathrm{reg}}(c)T,\qquad
C_{\mathrm{reg}}(c)=64C_{\mathrm{cmp}}+72C_{\mathrm{pr}}(c).
\label{reg:projected_fourth}
\end{equation}
\end{lemma}

\begin{proof}
The continuous fourth-moment bound \eqref{eq:R_fourth}, with the
$O(1/T)$ remainder in $N-\NO-S$, proves \eqref{reg:prime_fourth}.
Cover $[T/2,3T]$ by finitely many dyadic intervals. On each of them
keep the same cutoff $X=T^c$: its exponent relative to that interval
is $c+O_c(1/\log T)<1/2$, and the corresponding prime variance is
$O_c(1)$. This is why the assertion is for fixed $c<1/2$.
By minimality and $|z+w|^4\le8(|z|^4+|w|^4)$,
\[
\int|\widehat F_{T,c}-F_X|^4
\le\int|P_N-F_X|^4
\le8(C_{\mathrm{cmp}}+C_{\mathrm{pr}}(c))T.
\]
A second application of the same inequality gives
\eqref{reg:projected_fourth} with the stated constant.
\end{proof}

\subsection{Fourth moments on a clock-defined set of indices}
\begin{theorem}\label{reg:clock_theorem}
Fix $0<c<1/2$. For all sufficiently large $T$, every target with
$k\in\mathcal I_T^\tau=\{k:T\le\tau_k<2T\}$ is in the range of
$\widehat F_{T,c}$, so \eqref{reg:clock} defines the clock there.
Writing $\widehat u_{T,c,k}=k-\frac12-\widehat F_{T,c}(\gamma_k)$,
\begin{equation}
\left\langle|\widehat u_{T,c,k}|^4\right\rangle_{\mathcal I_T^\tau}
\ll_c1,\qquad
\left\langle|\widehat x_{T,c,k}-\gamma_k|^4\right\rangle_{\mathcal I_T^\tau}
\ll_c\lL(T)^{-4}.
\label{reg:fourth_clock}
\end{equation}
Consequently
\begin{equation}
\left\langle|\widehat x_{T,c,k}-\gamma_k|^2\right\rangle_{\mathcal I_T^\tau}
\ll_c\log^{-2}T,\qquad
\frac{\#\{k\in\mathcal I_T^\tau:
\lL(T)|\widehat x_{T,c,k}-\gamma_k|>R\}}
{\#\mathcal I_T^\tau}\ll_c R^{-4}\quad(R>0).
\label{reg:second_tail}
\end{equation}
For consecutive targets in its range,
\begin{equation}
\frac1{2\lL(T)}\le
\widehat x_{T,c,k+1}-\widehat x_{T,c,k}
\le\frac2{\lL(T)}.
\label{reg:gaps}
\end{equation}
The same moment estimates hold over $\mathcal I_T$.
\end{theorem}

\begin{proof}
Abbreviate $C_{\mathrm{reg}}(c)$ to $C_{\mathrm{reg}}$.
Equation~\eqref{reg:projected_fourth} supplies non-ordinates
$\alpha\in[T/2,3T/4]$ and $\beta\in[9T/4,5T/2]$ with
$|N(\alpha)-\widehat F_{T,c}(\alpha)|,
|N(\beta)-\widehat F_{T,c}(\beta)|\le(8C_{\mathrm{reg}})^{1/4}$.
On each of the two intervals, Markov's inequality leaves a set of
measure at least $T/8$ where this holds, so non-ordinates may be chosen.
Zero counting gives $N(T)-N(\alpha)\asymp T\log T$ and
$N(\beta)-N(2T)\asymp T\log T$.
Thus the target levels of $\mathcal I_T$ lie strictly between
$\widehat F_{T,c}(\alpha)$ and $\widehat F_{T,c}(\beta)$.
The same is true for $\mathcal I_T^\tau$, because
$|\tau_k-\gamma_k|=O(1)$ and their endpoint indices differ from those
of $\mathcal I_T$ by $O(\log T)$. All corresponding ordinates lie in
$(\alpha,\beta)$ for large $T$. This proves existence of the required
roots, using ordinates only in the proof.

For an absolutely continuous increasing function $P$, the exact identity
of Theorem~\ref{thm:identity_general} has the following endpoint form:
\begin{equation}
\begin{split}
\sum_{\alpha<\gamma_k\le\beta}
\Psi_2\bigl(k-\tfrac12-P(\gamma_k)\bigr)
&=\int_\alpha^\beta(N-P)^4P'\,dt\\
&\quad+\frac{(N(\beta)-P(\beta))^5
-(N(\alpha)-P(\alpha))^5}{5}.
\end{split}
\label{reg:endpoint_identity}
\end{equation}
Indeed $(N-P)^5/5$ has derivative $-(N-P)^4P'$ between ordinates;
at a jump of multiplicity $d$, splitting into $d$ unit jumps gives
the displayed sum. The chain rule holds for absolutely continuous $P$,
so the finitely many knots of a linear interpolant cause no extra terms.
Moreover $\Psi_2(u)=u^4+u^2/2+1/80\ge u^4$.

Apply \eqref{reg:endpoint_identity} with $P=\widehat F_{T,c}$.
The integral is at most $2\lL(T)C_{\mathrm{reg}}T$ and the endpoint
term is $O_c(1)$. Nonnegativity allows restriction of the sum to
either $\mathcal I_T$ or $\mathcal I_T^\tau$.
Their cardinalities are $\asymp T\lL(T)$, proving the first bound
of \eqref{reg:fourth_clock}. The lower slope gives
$|\widehat x_{T,c,k}-\gamma_k|\le
2|\widehat u_{T,c,k}|/\lL(T)$.
This proves the fourth-moment bound. Cauchy--Schwarz and Markov's
inequality give \eqref{reg:second_tail}. The two slope bounds give
\eqref{reg:gaps} directly.
\end{proof}

\subsection{Bounded transport with a fixed exponent}
\begin{theorem}\label{reg:transport_theorem}
Fix $0<c<1/4$ and $X=T^c$. For the clock just constructed,
\begin{equation}
\left\langle\left|
S_X(\gamma_k)-S_X(\widehat x_{T,c,k})
\right|^2\right\rangle_{\mathcal I_T^\tau}\ll_c1.
\label{reg:transport_bound}
\end{equation}
The same estimate holds over $\mathcal I_T$.
\end{theorem}

\begin{proof}
Let $d=|\widehat x_{T,c,k}-\gamma_k|$. On the interval joining the
two points, monotonicity of $N$ gives $|N-\widehat F_{T,c}|$ at least
$\mathsf a_T^{\mathrm{reg}}$ times the distance from $\widehat x_{T,c,k}$.
For $\widehat x_{T,c,k}>\gamma_k$, use $N(t)\ge k$;
for $\widehat x_{T,c,k}<\gamma_k$, use $N(t)\le k-1$ before $\gamma_k$.
These statements remain valid when ordinates repeat. Thus
\[
\frac{(\mathsf a_T^{\mathrm{reg}})^4d^5}{5}
\le\int_{\mathcal J_T^{\mathrm{reg}}}|N-\widehat F_{T,c}|^4\,dt,
\qquad
d\ll_c T^{1/5}\lL(T)^{-4/5}<T^{1/2}.
\]
The points lie in $[T-2T^{1/2},2T+2T^{1/2}]$ for either index set.
The gap lower bound \eqref{reg:gaps} supplies the multiplicity bound
directly: an interval of length $\mathfrak b_T=2\pi/\log(3T)$ contains
at most three points, since
$\mathfrak b_T<1/\lL(T)$ and three consecutive gaps span at least
$3/(2\lL(T))$. Thus $\tilde n\le3$ and
$\int\tilde n^2\le3\int\tilde n=3\mathfrak b_T\#I$
for either index set $I$. No near-pair estimate for zeros is needed
in this step.

Take a smooth cutoff $\varphi$ equal to $1$ on the localized interval,
supported in $[T/2,3T]$, with $|\varphi'|\ll1/T$, and set
$g=\mathcal M(\varphi S'_X)$. Lemma~\ref{lem:sampling_rho}, with
its bounded absolute constant in the sampling proof, yields
\[
\langle|g(\widehat x_{T,c,k})|^4\rangle\ll_c(1+\log X)^4.
\]
Here only fixed eighth derivative moments are needed. The restriction
$c<1/4$ ensures $X^4=o(T)$, so Montgomery--Vaughan applies on the
bounded dyadic covering of the cutoff support.
The maximal-function inequality and Cauchy--Schwarz now give
\[
\begin{split}
\langle|S_X(\gamma_k)-S_X(\widehat x_{T,c,k})|^2\rangle
&\le4\langle|\gamma_k-\widehat x_{T,c,k}|^4\rangle^{1/2}
\langle|g(\widehat x_{T,c,k})|^4\rangle^{1/2}\\
&\ll_c\frac{(1+\log X)^2}{\lL(T)^2}\ll_c1.
\end{split}
\]
\end{proof}

\subsection{Certified inexact minimization}
The construction need not solve \eqref{reg:projection} exactly.
For a prescribed $0\le\eta_{\mathrm{reg}}\le1$, let
$P_\eta=\mathsf p_{\mathbf v,T}^{\mathrm{reg}}\in\mathcal C_T^{\mathrm{reg}}$
satisfy
\begin{equation}
\mathscr E_{T,c}^{\mathrm{reg}}(\mathbf v)
\le\inf\mathscr E_{T,c}^{\mathrm{reg}}+\eta_{\mathrm{reg}}T.
\label{reg:tolerance}
\end{equation}
All previous arguments apply to this interpolant with
$C_{\mathrm{reg}}(c)$ replaced by $C_{\mathrm{reg}}(c)+8\eta_{\mathrm{reg}}$.
Consequently, uniformly in this tolerance range and for sufficiently
large $T$ depending only on $c$,
\begin{equation}
\begin{split}
\left\langle|\widehat x_{T,c,k}^{(\eta_{\mathrm{reg}})}-\gamma_k|^4
\right\rangle_{\mathcal I_T^\tau}
&\ll_c\frac{1+\eta_{\mathrm{reg}}}{\log^4T},\\
\left\langle|\widehat x_{T,c,k}^{(\eta_{\mathrm{reg}})}-\gamma_k|^2
\right\rangle_{\mathcal I_T^\tau}
&\ll_c\frac{1+\sqrt{\eta_{\mathrm{reg}}}}{\log^2T}.
\end{split}
\label{reg:inexact_clock}
\end{equation}
The roots and spacing bounds are obtained from the feasible nodal vector
by \eqref{reg:inverse}; they require no zero-dependent stopping rule.

Here is an explicit sufficient certificate for \eqref{reg:tolerance}.
Let $D\mathbf v=(v_{j+1}-v_j)_j$, let $\boldsymbol\mu_+,\boldsymbol\mu_-$
be nonnegative vectors, and abbreviate the objective to $E(\mathbf v)$.
Put $h=\mathsf h_T^{\mathrm{reg}}$, $J=|\mathcal J_T^{\mathrm{reg}}|=5T/2$,
$a=\mathsf a_T^{\mathrm{reg}}$, $b=\mathsf b_T^{\mathrm{reg}}$, and
\begin{equation}
\begin{split}
\mathscr L(\mathbf v,\boldsymbol\mu_+,\boldsymbol\mu_-)
&=E(\mathbf v)+\boldsymbol\mu_+^{\mathsf T}(D\mathbf v-bh\mathbf1)
+\boldsymbol\mu_-^{\mathsf T}(ah\mathbf1-D\mathbf v),\\
\mathbf r_{\mathrm{reg}}&=\nabla E(\mathbf v)+D^{\mathsf T}
(\boldsymbol\mu_+-\boldsymbol\mu_-),\qquad
A_{\mathrm{reg}}=\frac{h^2}{108J},\\
B_{\mathrm{reg}}&=\mathscr L(\mathbf v,\boldsymbol\mu_+,\boldsymbol\mu_-)
-\frac3{4^{4/3}}A_{\mathrm{reg}}^{-1/3}
\|\mathbf r_{\mathrm{reg}}\|_2^{4/3}.
\end{split}
\label{reg:dual_certificate}
\end{equation}
Then $B_{\mathrm{reg}}\le\inf E$. Therefore feasibility, nonnegative
multipliers, and $E(\mathbf v)-B_{\mathrm{reg}}\le\eta_{\mathrm{reg}}T$
certify \eqref{reg:tolerance}. For a numerical certificate all integrals,
gradients and final inequalities must be enclosed with directed rounding.

The objective and its gradient use only the known prime model. If
$\mathsf B_{T,i}^{\mathrm{reg}}$ is the nodal hat function, then
\[
\frac{\partial E}{\partial v_i}=4\int_{\mathcal J_T^{\mathrm{reg}}}
\bigl(\mathsf p_{\mathbf v,T}^{\mathrm{reg}}-F_X\bigr)^3
\mathsf B_{T,i}^{\mathrm{reg}}\,dt.
\]

To verify the lower bound, use the scalar identity
\[
(x+d)^4-x^4-4x^3d-\tfrac13d^4=6d^2(x+d/3)^2\ge0.
\]
The linear-interpolation mass matrix has smallest eigenvalue at least
$h/6$ by its explicit tridiagonal entries. Hence
\[
\int|\mathsf p_{\mathbf d,T}^{\mathrm{reg}}|^4
\ge J^{-1}\left(\int|\mathsf p_{\mathbf d,T}^{\mathrm{reg}}|^2\right)^2
\ge\frac{h^2}{36J}\|\mathbf d\|_2^4.
\]
It follows that $\mathscr L(\mathbf v+\mathbf d)\ge
\mathscr L(\mathbf v)+\mathbf r_{\mathrm{reg}}^{\mathsf T}\mathbf d
+A_{\mathrm{reg}}\|\mathbf d\|_2^4$.
Minimizing the last two terms over $\mathbf d$ gives
\eqref{reg:dual_certificate}. Weak duality proves the assertion.

\paragraph{Scope and constants.}
The construction itself has no unspecified slope, mesh, or stopping
parameter once $T,c,\eta_{\mathrm{reg}}$ have been fixed. The analytic
accuracy constants and the starting height still depend on the constants
in the fixed-order external moment bounds. Equations
\eqref{reg:comparator_constant} and \eqref{reg:projected_fourth} give a
dependency ledger, not a numerical finite-height accuracy certificate.
No numerical performance or historical priority is asserted here.
The bounds are blockwise in $T$ and do not by themselves define a single
global clock independent of the chosen height window. The bounded gap
range \eqref{reg:gaps} is imposed by the regularization; it is not a
claim about the fine spacing distribution of zeta zeros.
